\documentclass[11pt]{article}
\usepackage[margin=1in]{geometry}
\usepackage{booktabs}
\usepackage{longtable}
\usepackage{array}
\usepackage{hyperref}
\usepackage[T1]{fontenc}

\newcolumntype{F}{>{\ttfamily\footnotesize}p{0.32\textwidth}}
\newcolumntype{D}{>{\footnotesize}p{0.58\textwidth}}
\newcolumntype{C}{>{\ttfamily\footnotesize}p{0.08\textwidth}}
\newcolumntype{E}{>{\footnotesize}p{0.82\textwidth}}

\title{AFS (Air Facility System) Data Dictionary}
\author{}
\date{}

\begin{document}
\maketitle

\noindent Source: EPA, \emph{AFS Data Download} documentation (April 2014); ECHO AFS archive,
\url{https://echo.epa.gov/tools/data-downloads\#afs}.

\medskip
\noindent\textbf{Scope.} AFS is the pre-2014 Air Facility System; the ECHO download is \textbf{frozen as
of October 17, 2014}. Data are rolled up to the plant level. Each plant joins on \texttt{AFS\_ID} (also
called SCSC), a 10-character code (Census FIPS state + county + AFS plant number). AFS does not carry a
\texttt{REGISTRY\_ID}, so there is no built-in crosswalk to ICIS-Air's \texttt{PGM\_SYS\_ID}. Field names,
types, and code labels below follow EPA's AFS documentation; field names and the \emph{observed} code
values have been verified against the actual downloaded CSVs (discrepancies noted).

\section*{AFS\_FACILITIES.csv}
One row per plant: identification, location, classification, operating, and compliance status.
\begin{longtable}{F D}
\toprule
Field & Description \\ \midrule \endhead
PLANT\_ID & Numeric plant identifier \\
AFS\_ID & 10-character plant code (state FIPS + county FIPS + AFS plant number); also ``SCSC'' \\
PLANT\_NAME & Plant name \\
EPA\_REGION & EPA region (01--10) \\
PLANT\_STREET\_ADDRESS / PLANT\_CITY / PLANT\_COUNTY / STATE / STATE\_NUMBER / ZIP\_CODE & Location fields \\
PRIMARY\_SIC\_CODE / SECONDARY\_SIC\_CODE & Four-character Standard Industrial Classification codes \\
NAICS\_CODE & Six-character NAICS code \\
AFS\_GOV\_FACILITY\_CODE & Government facility indicator (see codes below) \\
FEDERALLY\_REPORTABLE & Y/N; ECHO-generated (see definition below) \\
EPA\_CLASSIFICATION\_CODE & Emissions classification (see codes below) \\
OPERATING\_STATUS & Operational condition (see codes below) \\
EPA\_COMPLIANCE\_STATUS & EPA compliance determination (see codes below) \\
CURRENT\_HPV & Current High Priority Violator status (see codes below) \\
LOCAL\_CONTROL\_REGION & Two-character local control region code (meanings vary by state) \\
STATE\_COMPLIANCE\_STATUS & State agency compliance determination (same code set as EPA\_COMPLIANCE\_STATUS) \\
\bottomrule
\end{longtable}

\subsection*{EPA\_CLASSIFICATION\_CODE (Alabama Power decision / 1993 EPA guidance)}
\begin{longtable}{C E}
\toprule
Code & Description \\ \midrule \endhead
A  & Actual or potential emissions above the applicable major-source thresholds \\
A1 & Actual or potential controlled emissions $>$100 tons/year \\
A2 & Actual emissions $<$100 tons/year, but potential uncontrolled emissions $>$100 tons/year \\
B  & Potential uncontrolled emissions $<$100 tons/year \\
C  & Unregulated pollutant, actual or potential controlled emissions $>$100 tons/year \\
ND & Major-source thresholds not defined \\
SM & Synthetic minor: potential emissions below all applicable major-source limits \\
UK & Unknown classification \\
\bottomrule
\end{longtable}
\noindent\emph{Correction:} \texttt{C} is \emph{not} ``Unknown'' (only \texttt{UK} is). Observed in the
data: A, A1, A2, B, C, ND, SM, UK.

\subsection*{OPERATING\_STATUS}
\begin{longtable}{C E}
\toprule
Code & Description \\ \midrule \endhead
O & Operating \\ C & Under Construction \\ P & Planned (applied for a construction permit) \\
T & Temporarily Closed \\ X & Permanently Closed \\ I & Seasonal \\
D & NESHAP Demolition \\ R & NESHAP Renovation \\ S & NESHAP Spraying \\ L & Landfill \\
\bottomrule
\end{longtable}
\noindent Observed in the data: O, C, P, T, X, I (the NESHAP/landfill statuses D, R, S, L are documented
but do not appear in this download).

\subsection*{EPA\_COMPLIANCE\_STATUS (worst case across EPA and state fields)}
\begin{longtable}{C E}
\toprule
Code & Description \\ \midrule \endhead
0 & Unknown compliance status \\
1 & In Violation -- No Schedule \\
2 & In Compliance -- Source Test \\
3 & In Compliance -- Inspection \\
4 & In Compliance -- Certification \\
5 & Meeting Compliance Schedule \\
6 & In Violation -- Not Meeting Schedule \\
7 & In Violation -- Unknown With Regard To Schedule \\
8 & No Applicable State Regulation \\
9 & In Compliance -- Shut Down \\
A & Unknown With Regard To Procedural Compliance \\
B & In Violation With Regard To Both Emissions And Procedural Compliance \\
C & In Compliance With Procedural Requirements \\
D & HPV Violation (auto-generated) \\
E & FRV Violation (auto-generated) \\
F & HPV On Schedule (auto-generated) \\
G & FRV On Schedule (auto-generated) \\
H & In Compliance (auto-generated) \\
M & In Compliance -- CEMs \\
P & Present, See Other Program(s) \\
U & Unknown By Evaluation Calculation \\
W & In Violation With Regard To Procedural Compliance \\
Y & Unknown With Regard To Both Emissions And Procedural Compliance \\
\bottomrule
\end{longtable}
\noindent Codes A, B, C, P, U, W, Y appear in the data and are defined in EPA's AFS documentation; a
prior draft listed only 0--9, D, E, F, G, H, M.

\subsection*{CURRENT\_HPV}
\begin{longtable}{C E}
\toprule
Code & Description \\ \midrule \endhead
S & Violation Unaddressed; State/Local has lead enforcement \\
T & Violation Addressed; State has lead enforcement \\
E & Violation Unaddressed; EPA has lead enforcement \\
F & Violation Addressed; EPA has lead enforcement \\
B & Violation Unaddressed; EPA and State share lead enforcement \\
C & Violation Addressed; EPA and State share lead enforcement \\
X & Violation Unaddressed; enforcement lead unassigned \\
\bottomrule
\end{longtable}
\noindent Observed in the data: C, E, F, S, T (B and X are documented but do not appear).

\subsection*{AFS\_GOV\_FACILITY\_CODE}
\texttt{0} Privately owned/operated \textbar{} \texttt{1} Federal \textbar{} \texttt{2} State \textbar{}
\texttt{3} County \textbar{} \texttt{4} Municipality \textbar{} \texttt{5} District \textbar{}
\texttt{6} Tribe. (All 0--6 observed.)

\subsection*{FEDERALLY\_REPORTABLE}
ECHO-generated Y/N. Per EPA documentation, ``Y'' if \texttt{EPA\_CLASSIFICATION\_CODE} $\in$
\{A, A1, A2, SM\}, OR (\texttt{AIR\_PROGRAM\_CODE} $\in$ \{8, 9\} and \texttt{EPA\_COMPLIANCE\_STATUS}
$\neq$ 8).

\section*{AIR\_PROGRAM.csv}
One row per plant--program (with pollutant-level detail). Joins to facilities on \texttt{AFS\_ID}/\texttt{PLANT\_ID}.
\begin{longtable}{F D}
\toprule
Field & Description \\ \midrule \endhead
AFS\_ID / PLANT\_ID & Plant identifiers \\
AIR\_PROGRAM\_CODE & One-character regulatory air program (see codes below) \\
AIR\_PROGRAM\_STATUS & Operating status within the program (same codes as OPERATING\_STATUS) \\
EPA\_CLASSIFICATION\_CODE & Emissions classification at the air-program level \\
EPA\_COMPLIANCE\_STATUS & Compliance status at the air-program level \\
AIR\_PROGRAM\_CODE\_SUBPARTS & Applicable subparts, space-delimited \\
POLLUTANT\_CODE & Pollutant code at the air-program level \\
CHEMICAL\_ABSTRACT\_SERVICE\_NMBR & CAS number for the pollutant, if any \\
POLLUTANT\_CLASSIFICATION & Emissions classification at the pollutant level \\
POLLUTANT\_COMPLIANCE\_STATUS & Compliance status at the pollutant level \\
\bottomrule
\end{longtable}

\subsection*{AIR\_PROGRAM\_CODE (per EPA AFS documentation)}
\texttt{0} SIP \textbar{} \texttt{1} SIP under federal jurisdiction (FIP) \textbar{}
\texttt{3} Non-Federally Reportable \textbar{} \texttt{4} CFC Tracking \textbar{} \texttt{6} PSD \textbar{}
\texttt{7} NSR \textbar{} \texttt{8} NESHAP (Part 61) \textbar{} \texttt{9} NSPS \textbar{}
\texttt{A} Acid Precipitation \textbar{} \texttt{F} FESOP (non-Title V) \textbar{}
\texttt{I} Native American \textbar{} \texttt{M} MACT (Part 63) \textbar{}
\texttt{T} TIP (Tribal Implementation Plan) \textbar{} \texttt{V} Title V.

\noindent\emph{Data note:} the \texttt{AIR\_PROGRAM\_CODE} column also contains the values \texttt{G} and
\texttt{R}, which are \emph{not} in EPA's documented program-code list above; their meaning is not
documented in the AFS download materials.

\section*{AFS\_ACTIONS.csv}
Compliance-monitoring and enforcement events, rolled up to plant level (1978--present).
\begin{longtable}{F D}
\toprule
Field & Description \\ \midrule \endhead
AFS\_ID / PLANT\_ID & Plant identifiers \\
ANU1 & Action number; uniquely identifies an action within a plant \\
NATIONAL\_ACTION\_TYPE & Two-character action code (inspection / enforcement). See note below \\
NATIONAL\_ACTION\_DESC & Text description of the action type \\
DATE\_ACHIEVED & Date the action was completed (YYYYMMDD) \\
ALL\_AIR\_PROGRAM\_CODES & All air programs associated with the action, space-delimited \\
PENALTY\_AMOUNT & Civil penalty assessed or agreed, in dollars \\
RESULT\_CODE & Result of stack tests / Title V reviews / actions \\
POLLUTANT\_CODE & Pollutant associated with the action \\
ALL\_VIOLATING\_POLL\_CODES & Pollutant(s) in violation, space-delimited \\
ALL\_VIOLATION\_TYPE\_CODES & Violation type code(s) \\
KEY\_ACTION\_NUMBERS & Links to violation / FCE pathways (up to ten) \\
REGIONAL\_DATA\_ELEMENT\_8 & Region-specific data element \\
CREATION\_DATE / DATE\_RECORD\_IS\_UPDATED & Record metadata dates \\
\bottomrule
\end{longtable}
\noindent\textbf{NATIONAL\_ACTION\_TYPE} translates region-specific codes to EPA national activity codes;
the lead agency is indicated within each code's description. Per EPA, the most common codes are
\texttt{FF, FS, FE, FZ, 1A, 5C} (full inspections); \texttt{EM, EO, ES, EX, PC, PO, PP, PR, PS, PX}
(partial inspections); and \texttt{1B, 2D, 6B, 7A, 7E, 7F, 8A, 8C, 9A} (formal enforcement). The download
contains \textbf{110 distinct} \texttt{NATIONAL\_ACTION\_TYPE} values and over 130 \texttt{RESULT\_CODE}
values; the full code lists are given in EPA's AFS documentation, and each row also carries its own
\texttt{NATIONAL\_ACTION\_DESC}.

\section*{AFS\_AIR\_PRG\_HIST\_COMPLIANCE.csv}
Quarterly compliance status per plant--program (FY2007--present). One row per plant--program--quarter.
\begin{longtable}{F D}
\toprule
Field & Description \\ \midrule \endhead
AFS\_ID & Plant identifier \\
AIR\_PROGRAM\_CODE & Regulatory program (same codes as AIR\_PROGRAM) \\
HISTORICAL\_COMPLIANCE\_DATE & Quarter, YYQQ (Q1 = Jan--Mar, \dots, Q4 = Oct--Dec) \\
HISTORICAL\_COMPLIANCE\_STATUS & Compliance status for the quarter (same codes as EPA\_COMPLIANCE\_STATUS) \\
\bottomrule
\end{longtable}

\section*{AFS\_HPV\_HISTORY.csv}
Lifecycle of High Priority Violator designations. One row per HPV episode.
\begin{longtable}{F D}
\toprule
Field & Description \\ \midrule \endhead
AFS\_ID & Plant identifier \\
HPV\_DAYZERO\_TYPE & Lead-agency code at day zero (carried with HPV\_DAYZERO\_DESC) \\
HPV\_DAYZERO\_DESC & Text description of the day-zero type \\
HPV\_DAYZERO\_DATE & Date the plant entered HPV status \\
HPV\_RESOLVED\_TYPE & Resolution-action code (blank if unresolved; carried with HPV\_RESOLVED\_DESC) \\
HPV\_RESOLVED\_DESC & Text description of the resolution type \\
HPV\_RESOLVED\_DATE & Date the HPV was resolved \\
\bottomrule
\end{longtable}
\noindent The \texttt{*\_TYPE} code meanings are carried in the adjacent \texttt{*\_DESC} columns in the
data itself, and are listed in EPA's AFS documentation.

\end{document}
